\chapter{DX-042 Shared UI Scene IR And Bijou Render Target}\label{record-138}

\noindent{\footnotesize\color{BookMuted}\path{docs/design/DX-042-shared-ui-scene-ir-and-bijou-render-target.md}\par}\smallskip

\begin{RecordMeta}
title: DX-042 Shared UI Scene IR And Bijou Render Target
legend: DX
lane: cool-ideas
priority: medium
github_issue: 302
status: proposed
keywords:
  - developer-experience
  - ir
  - geordi
  - wesley
  - surface
  - provenance

\end{RecordMeta}

Legend: \href{https://github.com/flyingrobots/bijou/blob/1853ad902f6bf219326409e7f743a5a1d6517dbc/docs/legends/DX-developer-experience.md}{DX - Developer Experience}

\section{Linked Work}\label{record-138-linked-work}

\begin{itemize}
\tightlist
\item
  User story: \href{https://github.com/flyingrobots/bijou/issues/302}{\#302}
\item
  Roadmap candidate: Runtime Graph And Scene IR
\item
  This document is a planning artifact only. It records the shared IR boundary for the \#302 idea; it does not claim implementation of the IR, a compiler, or a renderer integration in this PR.
\end{itemize}

\section{Decision Summary}\label{record-138-decision-summary}

Bijou \texttt{S\allowbreak{}u\allowbreak{}r\allowbreak{}f\allowbreak{}a\allowbreak{}c\allowbreak{}e} should be treated as a deterministic terminal render target for a portable UI scene IR, not as the portable IR itself.

The shared IR should sit above Bijou and Geordi:

\begin{CodeBlock}
GraphQL SDL / authored source
  -> Wesley semantic graph and domain law
    -> ui-scene-ir/1
      -> bijou-layout/1
        -> Bijou Surface
          -> terminal cells

      -> geordi-ir/1
        -> Geordi renderer
          -> pixels and render receipts
\end{CodeBlock}

This keeps the final renderer substrates honest. Geordi owns pixel-oriented render proof. Bijou owns terminal cell rendering, input, focus, runtime interaction, and agent-legible surfaces. The shared IR owns portable scene structure, semantic identity, bindings, localization references, theme token references, actions, and source maps.

\section{Sponsored Human}\label{record-138-sponsored-human}

A builder wants to describe a UI scene once and render it through either Geordi or Bijou so that browser/native artifacts and terminal artifacts stay aligned without manually maintaining two different scene descriptions.

\section{Sponsored Agent}\label{record-138-sponsored-agent}

An agent wants a source-mapped, inspectable scene contract so it can audit which localized strings, theme tokens, bindings, actions, and rendered cells came from the same authored scene.

\section{Hill}\label{record-138-hill}

A scene authored as a semantic UI contract can lower to both a Geordi render artifact and a Bijou terminal surface, with structural receipts proving that node identity, text keys, token references, bindings, and actions stayed in sync across both targets.

\section{Core Rule}\label{record-138-core-rule}

The portable IR describes scene intent.

Bijou \texttt{S\allowbreak{}u\allowbreak{}r\allowbreak{}f\allowbreak{}a\allowbreak{}c\allowbreak{}e} describes rendered terminal output.

Those are different layers:

\begin{CodeBlock}
ui-scene-ir/1
  describes: nodes, text keys, token refs, bindings, actions, source maps

bijou-layout/1
  describes: terminal layout, focus graph, keymap, viewport regions

Bijou Surface
  describes: final cell grid with glyphs, colors, modifiers, opacity, masks
\end{CodeBlock}

Bijou may render a shared scene IR. Geordi may render the same shared scene IR. Neither system should pretend that terminal cells and pixels are the same final substrate.

\section{Playback Questions}\label{record-138-playback-questions}

\begin{enumerate}
\def\labelenumi{\arabic{enumi}.}
\tightlist
\item
  Which authored scene node produced this rendered cell or pixel region?
\item
  Which localization key produced this text?
\item
  Which theme token produced this foreground or background color?
\item
  Which action or key binding is attached to this node?
\item
  Which renderer-specific target profile lowered this scene?
\item
  Which target-specific claims are proven, and which are explicit nonclaims?
\item
  Can the Geordi and Bijou outputs prove structural parity without claiming pixel parity?
\item
  Can compile-time checks catch missing localization keys and raw color usage before runtime?
\end{enumerate}

\section{Boundary}\label{record-138-boundary}

This design is about a portable scene contract and renderer-specific lowering.

It is not about:

\begin{itemize}
\tightlist
\item
  replacing Bijou's runtime loop
\item
  replacing Bijou \texttt{S\allowbreak{}u\allowbreak{}r\allowbreak{}f\allowbreak{}a\allowbreak{}c\allowbreak{}e}
\item
  making terminal cells behave like pixels
\item
  making Geordi own terminal input or focus
\item
  making Wesley a renderer
\item
  forcing every existing Bijou block to be authored in GraphQL immediately
\item
  claiming pixel-perfect parity between terminal and browser render targets
\end{itemize}

\section{Role Model}\label{record-138-role-model}

\subsection{Wesley}\label{record-138-wesley}

Wesley owns semantic compilation and domain law.

For this design, Wesley can:

\begin{itemize}
\tightlist
\item
  parse GraphQL SDL or another schema-backed authoring source
\item
  validate schema shape
\item
  generate TypeScript and Rust types
\item
  validate localization requirements
\item
  validate style token requirements
\item
  compute schema hashes
\item
  produce or validate \texttt{u\allowbreak{}i\allowbreak{}-\allowbreak{}s\allowbreak{}c\allowbreak{}e\allowbreak{}n\allowbreak{}e\allowbreak{}-\allowbreak{}i\allowbreak{}r\allowbreak{}/\allowbreak{}1}
\end{itemize}

Wesley should not render terminal cells or pixels.

\subsection{Shared UI Scene IR}\label{record-138-shared-ui-scene-ir}

\texttt{u\allowbreak{}i\allowbreak{}-\allowbreak{}s\allowbreak{}c\allowbreak{}e\allowbreak{}n\allowbreak{}e\allowbreak{}-\allowbreak{}i\allowbreak{}r\allowbreak{}/\allowbreak{}1} owns portable scene structure.

It should carry enough information for both Bijou and Geordi to lower the same scene while preserving source identity.

It should know about:

\begin{itemize}
\tightlist
\item
  scene node IDs
\item
  node kinds
\item
  semantic roles
\item
  layout intent
\item
  text and localization references
\item
  theme token references
\item
  bindings
\item
  actions
\item
  source maps
\item
  target profile requirements
\end{itemize}

It should not know about:

\begin{itemize}
\tightlist
\item
  terminal escape sequences
\item
  renderer dirty-region internals
\item
  final RGB cell caches
\item
  subpixel antialiasing details
\item
  host-specific input event loops
\end{itemize}

\subsection{Geordi}\label{record-138-geordi}

Geordi owns pixel-oriented render artifacts and receipts.

For this design, Geordi lowers \texttt{u\allowbreak{}i\allowbreak{}-\allowbreak{}s\allowbreak{}c\allowbreak{}e\allowbreak{}n\allowbreak{}e\allowbreak{}-\allowbreak{}i\allowbreak{}r\allowbreak{}/\allowbreak{}1} to \texttt{g\allowbreak{}e\allowbreak{}o\allowbreak{}r\allowbreak{}d\allowbreak{}i\allowbreak{}-\allowbreak{}i\allowbreak{}r\allowbreak{}/\allowbreak{}1}, then renders or verifies the pixel artifact under a declared target profile.

Geordi claims:

\begin{itemize}
\tightlist
\item
  explicit pixel geometry
\item
  font and asset requirements
\item
  visual receipt hashes
\item
  target profile facts
\item
  browser/native renderer parity where supported
\end{itemize}

Geordi does not claim:

\begin{itemize}
\tightlist
\item
  terminal focus semantics
\item
  terminal key routing
\item
  cell-level glyph approximations
\item
  Bijou runtime state updates
\end{itemize}

\subsection{Bijou}\label{record-138-bijou}

Bijou owns terminal runtime behavior and terminal render targets.

For this design, Bijou lowers \texttt{u\allowbreak{}i\allowbreak{}-\allowbreak{}s\allowbreak{}c\allowbreak{}e\allowbreak{}n\allowbreak{}e\allowbreak{}-\allowbreak{}i\allowbreak{}r\allowbreak{}/\allowbreak{}1} to \texttt{b\allowbreak{}i\allowbreak{}j\allowbreak{}o\allowbreak{}u\allowbreak{}-\allowbreak{}l\allowbreak{}a\allowbreak{}y\allowbreak{}o\allowbreak{}u\allowbreak{}t\allowbreak{}/\allowbreak{}1}, then renders that layout into a \texttt{S\allowbreak{}u\allowbreak{}r\allowbreak{}f\allowbreak{}a\allowbreak{}c\allowbreak{}e}.

Bijou claims:

\begin{itemize}
\tightlist
\item
  deterministic terminal cell output
\item
  focus and input routing
\item
  theme token resolution
\item
  localization lowering
\item
  cell-level provenance
\item
  surface receipts
\item
  agent-readable inspection through structured surfaces
\end{itemize}

Bijou does not claim:

\begin{itemize}
\tightlist
\item
  browser pixel parity
\item
  native canvas fidelity
\item
  subpixel antialiasing
\item
  font-pack proof identical to Geordi
\end{itemize}

\section{Relationship Diagram}\label{record-138-relationship-diagram}

\begin{center}
\resizebox{\linewidth}{!}{%
\begin{tikzpicture}[booknode/.style={draw=BookHeading,line width=.6pt,fill=BookPanel,rounded corners=1.5pt,align=center,text width=29mm,minimum height=10mm,inner sep=5pt,font=\small},bookedge/.style={-{Latex[length=2mm]},draw=BookHeading,line width=.65pt},every edge quotes/.style={font=\scriptsize,fill=white,inner sep=2pt,text=BookMuted}]
\node[booknode] (SDL) at (0,0) {Graph\allowbreak{}QL SDL or authored source};
\node[booknode] (W) at (0,-1.6) {Wesley semantic compiler};
\node[booknode] (IR) at (0,-3.2) {ui-scene-ir/1};
\node[booknode] (BIR) at (-3.5,-4.9) {bijou-layout/1};
\node[booknode] (Surface) at (-3.5,-6.5) {Bijou Surface};
\node[booknode] (Terminal) at (-5,-8.2) {Terminal cells};
\node[booknode] (GIR) at (3.5,-4.9) {geordi-ir/1};
\node[booknode] (Pixels) at (1.8,-6.8) {Pixels};
\node[booknode] (GR) at (5.2,-6.8) {Geordi receipt};
\node[booknode] (BR) at (-1.5,-8.2) {Bijou surface receipt};
\draw[bookedge,] (SDL) to (W);
\draw[bookedge,] (W) to (IR);
\draw[bookedge,] (IR) to (BIR);
\draw[bookedge,] (BIR) to (Surface);
\draw[bookedge,] (Surface) to (Terminal);
\draw[bookedge,] (Surface) to (BR);
\draw[bookedge,] (IR) to (GIR);
\draw[bookedge,] (GIR) to (Pixels);
\draw[bookedge,] (GIR) to (GR);
\end{tikzpicture}}
\end{center}

\section{Target Profile Diagram}\label{record-138-target-profile-diagram}

\begin{CodeBlock}
ui-scene-ir/1
  targetProfiles:
    - kind: geordi-pixel
      width: 1280
      height: 720
      colorProfile: truecolor
      claims:
        - pixel geometry
        - font metrics
        - asset hashes

    - kind: bijou-cell
      cols: 120
      rows: 40
      glyphProfile: ansi-unicode
      claims:
        - structural hierarchy
        - cell occupancy
        - token provenance
        - i18n provenance
\end{CodeBlock}

The target profile makes renderer differences explicit. A Bijou terminal render can claim structural parity and cell-level determinism without claiming that it matches Geordi pixels exactly.

\section{Proposed IR Shape}\label{record-138-proposed-ir-shape}

The shared scene IR should be plain data. The first draft can stay intentionally small:

\begin{CodeBlock}
type UiSceneIr = {
  irVersion: 'ui-scene-ir/1';
  id: string;
  sourceHash: string;
  rootNodeId: string;
  nodes: UiNode[];
  bindings: UiBinding[];
  actions: UiAction[];
  tokenUses: UiTokenUse[];
  i18nUses: UiI18nUse[];
  sourceMap: UiSourceMapEntry[];
  targetProfiles: UiTargetProfile[];
};
\end{CodeBlock}

\subsection{Nodes}\label{record-138-nodes}

\begin{CodeBlock}
type UiNode = {
  id: string;
  kind: 'box' | 'text' | 'image' | 'group' | 'list' | 'table' | 'custom';
  role?: string;
  parentId?: string;
  children?: string[];
  layout?: UiLayoutIntent;
  text?: UiTextRef;
  style?: UiStyleRef;
  actions?: string[];
  metadata?: Record<string, unknown>;
};
\end{CodeBlock}

Nodes are semantic scene members. They are not terminal cells and they are not pixel rectangles, even when a renderer later lowers them into cells or rectangles.

\subsection{Text}\label{record-138-text}

\begin{CodeBlock}
type UiTextRef =
  | { kind: 'literal'; value: string }
  | { kind: 'i18n'; key: string; fallback?: string };
\end{CodeBlock}

Literal text should be allowed for fixture work and generated artifacts, but product surfaces should prefer \texttt{i\allowbreak{}1\allowbreak{}8\allowbreak{}n} references when they are user-facing.

\subsection{Style}\label{record-138-style}

\begin{CodeBlock}
type UiStyleRef = {
  fg?: { token: string };
  bg?: { token: string };
  border?: { token: string };
  modifiers?: string[];
};
\end{CodeBlock}

The IR should carry token references, not only resolved colors. Renderers can resolve tokens later against dark or light themes. Lower modes can preserve the token name when color is unavailable.

\subsection{Bindings}\label{record-138-bindings}

\begin{CodeBlock}
type UiBinding = {
  id: string;
  targetNodeId: string;
  targetProperty: string;
  source: {
    kind: 'state' | 'query' | 'computed';
    path: string;
  };
  when?: string;
};
\end{CodeBlock}

Bindings are data, not hidden callbacks. They make dynamic surfaces auditable and compile-time-checkable.

\subsection{Actions}\label{record-138-actions}

\begin{CodeBlock}
type UiAction = {
  id: string;
  label?: UiTextRef;
  command: string;
  keybindings?: string[];
  targetNodeId?: string;
};
\end{CodeBlock}

Actions let Bijou preserve focus and input contracts while Geordi can still render the same action affordances as visual artifacts.

\section{Bijou Lowering}\label{record-138-bijou-lowering}

Bijou lowering turns portable intent into a terminal-specific layout contract:

\begin{CodeBlock}
type BijouLayoutIr = {
  irVersion: 'bijou-layout/1';
  sceneHash: string;
  nodes: LayoutNode[];
  focusGraph: FocusGraph;
  keymap: KeyBinding[];
  viewportRegions: ViewportRegion[];
  themeRequirements: TokenRequirement[];
  localeRequirements: I18nRequirement[];
  sourceMap: UiToCellSourceMap[];
};
\end{CodeBlock}

The final render target remains \texttt{S\allowbreak{}u\allowbreak{}r\allowbreak{}f\allowbreak{}a\allowbreak{}c\allowbreak{}e}.

\begin{CodeBlock}
ui-scene-ir/1
  -> bijou-layout/1
    -> Surface
      -> terminal
\end{CodeBlock}

The \texttt{S\allowbreak{}u\allowbreak{}r\allowbreak{}f\allowbreak{}a\allowbreak{}c\allowbreak{}e} can then be inspected, hashed, diffed, serialized for MCP, or recorded in a renderer playback artifact.

\section{Geordi Lowering}\label{record-138-geordi-lowering}

Geordi lowering turns portable intent into a pixel-oriented render contract:

\begin{CodeBlock}
ui-scene-ir/1
  -> geordi-ir/1
    -> browser/native renderer
      -> pixel buffer
      -> receipt
\end{CodeBlock}

Geordi can preserve the same node IDs, token references, localization references, and source maps while resolving geometry and render-profile facts for pixel output.

\section{Cross-Renderer Structural Receipt}\label{record-138-cross-renderer-structural-receipt}

The first proof should not compare pixels to cells. It should compare shared structure.

\begin{CodeBlock}
type SharedSceneReceipt = {
  sceneHash: string;
  sourceHash: string;
  nodeIds: string[];
  i18nKeys: string[];
  tokenRefs: string[];
  actionIds: string[];
  bindingIds: string[];
  outputs: {
    geordi?: {
      irHash: string;
      receiptHash: string;
    };
    bijou?: {
      layoutHash: string;
      surfaceHash: string;
    };
  };
};
\end{CodeBlock}

This answers the important question:

\begin{CodeBlock}
Did both targets render the same authored scene contract?
\end{CodeBlock}

It does not falsely answer:

\begin{CodeBlock}
Are terminal cells pixel-perfect replicas of browser pixels?
\end{CodeBlock}

\section{Debug And Provenance Files}\label{record-138-debug-and-provenance-files}

This design connects directly to the proposed stable debug logs:

\begin{CodeBlock}
<repo>/.bijou/debug.i18n.jsonl
<repo>/.bijou/debug.style.jsonl
\end{CodeBlock}

The logs should record first appearance or provenance changes, not every frame.

Example i18n line:

\begin{CodeBlock}
{"scene":"docs.home","node":"hero.title","key":"dogfood.home.title","locale":"en","status":"resolved"}
\end{CodeBlock}

Example style line:

\begin{CodeBlock}
{"scene":"docs.home","node":"hero.title","slot":"fg","token":"semantic.primary","mode":"dark","status":"resolved"}
\end{CodeBlock}

The eventual goal is to catch these issues earlier through IR validation:

\begin{itemize}
\tightlist
\item
  missing localization key
\item
  unsupported locale
\item
  raw color in product surface
\item
  unresolved token
\item
  unsafe token pair
\item
  binding target that does not exist
\item
  action without a command
\end{itemize}

Runtime logs remain useful as proof of what actually appeared.

\section{Authoring Model}\label{record-138-authoring-model}

GraphQL SDL remains the strongest candidate authoring source because Wesley and Geordi already use schema-backed compilation patterns.

The authoring source may look like this in spirit:

\begin{CodeBlock}
type DocsHomeScene
  @ui_scene(id: "docs.home")
  @ui_target(kind: "bijou-cell", cols: 120, rows: 40)
  @ui_target(kind: "geordi-pixel", width: 1280, height: 720) {
  title: String!
    @ui_node(id: "hero.title", kind: "text")
    @ui_text(key: "dogfood.home.title")
    @ui_style(fg: "semantic.primary")

  openAction: String
    @ui_action(id: "docs.open", command: "docs.open", keybindings: ["Enter"])
}
\end{CodeBlock}

This is not final syntax. The important contract is that scene structure, text, style, binding, and action identity are compile-time data.

\section{Lower Modes}\label{record-138-lower-modes}

The shared IR should make lower modes more deterministic.

For example, Bijou could render diagnostic modes:

\begin{CodeBlock}
normal:
  Start Here

i18n lowering:
  dogfood.nav.startHere

token lowering:
  semantic.primary on surface.primary.bg

binding lowering:
  docs.currentPage.title <- route.currentPage.title
\end{CodeBlock}

These modes are easier to test than screenshots and more useful to agents than raw escape sequences.

\section{Scope}\label{record-138-scope}

The first implementation slice should include:

\begin{itemize}
\tightlist
\item
  a design document for \texttt{u\allowbreak{}i\allowbreak{}-\allowbreak{}s\allowbreak{}c\allowbreak{}e\allowbreak{}n\allowbreak{}e\allowbreak{}-\allowbreak{}i\allowbreak{}r\allowbreak{}/\allowbreak{}1}
\item
  one tiny fixture scene
\item
  a structural receipt format
\item
  a Bijou lowering proof to \texttt{S\allowbreak{}u\allowbreak{}r\allowbreak{}f\allowbreak{}a\allowbreak{}c\allowbreak{}e}
\item
  a Geordi lowering proof or adapter stub if Geordi integration is not yet ready inside this repository
\item
  tests proving node IDs, token refs, i18n keys, actions, and bindings survive lowering
\end{itemize}

\section{Non-Goals}\label{record-138-non-goals}

The first slice should not:

\begin{itemize}
\tightlist
\item
  migrate DOGFOOD wholesale to GraphQL-authored scenes
\item
  replace existing Bijou block APIs
\item
  implement a full Geordi renderer inside Bijou
\item
  require Geordi as a runtime dependency of every Bijou app
\item
  claim pixel/cell parity
\item
  solve every binding expression shape
\end{itemize}

\section{Acceptance Criteria}\label{record-138-acceptance-criteria}

\begin{itemize}
\tightlist
\item
  A \texttt{u\allowbreak{}i\allowbreak{}-\allowbreak{}s\allowbreak{}c\allowbreak{}e\allowbreak{}n\allowbreak{}e\allowbreak{}-\allowbreak{}i\allowbreak{}r\allowbreak{}/\allowbreak{}1} fixture describes at least one text node, one token reference, one i18n key, and one action.
\item
  Bijou can lower the fixture into a terminal render target without losing node identity.
\item
  The resulting \texttt{S\allowbreak{}u\allowbreak{}r\allowbreak{}f\allowbreak{}a\allowbreak{}c\allowbreak{}e} can report a structural receipt or source-map facts tied back to the fixture.
\item
  The design names the distinction between shared IR, Bijou layout IR, and Bijou \texttt{S\allowbreak{}u\allowbreak{}r\allowbreak{}f\allowbreak{}a\allowbreak{}c\allowbreak{}e}.
\item
  The design names Geordi as the pixel-proof target and Bijou as the terminal-surface target.
\item
  The validation plan proves structure, not screenshot similarity.
\end{itemize}

\section{Slice Plan}\label{record-138-slice-plan}

\begin{enumerate}
\def\labelenumi{\arabic{enumi}.}
\tightlist
\item
  Document the shared IR contract and renderer boundaries.
\item
  Add a minimal \texttt{u\allowbreak{}i\allowbreak{}-\allowbreak{}s\allowbreak{}c\allowbreak{}e\allowbreak{}n\allowbreak{}e\allowbreak{}-\allowbreak{}i\allowbreak{}r\allowbreak{}/\allowbreak{}1} fixture package or test helper.
\item
  Add Bijou lowering from the fixture to a simple \texttt{S\allowbreak{}u\allowbreak{}r\allowbreak{}f\allowbreak{}a\allowbreak{}c\allowbreak{}e}.
\item
  Add structural receipt generation for the fixture.
\item
  Add token and i18n provenance checks against the receipt.
\item
  Add a Geordi integration spike that either lowers to \texttt{g\allowbreak{}e\allowbreak{}o\allowbreak{}r\allowbreak{}d\allowbreak{}i\allowbreak{}-\allowbreak{}i\allowbreak{}r\allowbreak{}/\allowbreak{}1} or records the adapter boundary needed to do so.
\end{enumerate}

\section{Validation Plan}\label{record-138-validation-plan}

Focused proof should use deterministic fixture tests:

\begin{CodeBlock}
npm test -- --run shared-ui-scene-ir
npm run typecheck:test
npm run lint
\end{CodeBlock}

Future integration proof can add:

\begin{CodeBlock}
npm run dogfood:i18n:check
npm run dogfood:i18n:debt
npm run verify:interactive-examples
\end{CodeBlock}

The key invariant is that a passing test must prove software behavior or inspectable output. A design document alone is not completion proof for the implementation goalpost.

\section{Risks}\label{record-138-risks}

\subsection{Risk: The Shared IR Becomes A Renderer}\label{record-138-risk-the-shared-ir-becomes-a-renderer}

If \texttt{u\allowbreak{}i\allowbreak{}-\allowbreak{}s\allowbreak{}c\allowbreak{}e\allowbreak{}n\allowbreak{}e\allowbreak{}-\allowbreak{}i\allowbreak{}r\allowbreak{}/\allowbreak{}1} starts carrying renderer internals, it will become too expensive to author and too vague to validate.

Mitigation: keep final render details in \texttt{g\allowbreak{}e\allowbreak{}o\allowbreak{}r\allowbreak{}d\allowbreak{}i\allowbreak{}-\allowbreak{}i\allowbreak{}r\allowbreak{}/\allowbreak{}1}, \texttt{b\allowbreak{}i\allowbreak{}j\allowbreak{}o\allowbreak{}u\allowbreak{}-\allowbreak{}l\allowbreak{}a\allowbreak{}y\allowbreak{}o\allowbreak{}u\allowbreak{}t\allowbreak{}/\allowbreak{}1}, and \texttt{S\allowbreak{}u\allowbreak{}r\allowbreak{}f\allowbreak{}a\allowbreak{}c\allowbreak{}e}.

\subsection{Risk: Bijou Loses Terminal-Native Behavior}\label{record-138-risk-bijou-loses-terminal-native-behavior}

If the IR treats keyboard routing, focus, and terminal viewports as afterthought metadata, Bijou scenes will look portable but feel broken.

Mitigation: make Bijou target profiles carry focus, keymap, and viewport requirements explicitly.

\subsection{Risk: Geordi And Bijou Drift Anyway}\label{record-138-risk-geordi-and-bijou-drift-anyway}

If both lowerings are implemented independently without receipts, drift will come back.

Mitigation: structural receipts compare node IDs, text keys, token refs, bindings, actions, and source hashes.

\subsection{Risk: GraphQL Syntax Gets Overloaded}\label{record-138-risk-graphql-syntax-gets-overloaded}

GraphQL SDL can describe contracts well, but directives can become JSON string dumps if the schema is not designed carefully.

Mitigation: use Wesley domain law to validate directive shape and prefer typed schema contracts over ad hoc JSON blobs.

\section{Open Questions}\label{record-138-open-questions}

\begin{itemize}
\tightlist
\item
  Should \texttt{u\allowbreak{}i\allowbreak{}-\allowbreak{}s\allowbreak{}c\allowbreak{}e\allowbreak{}n\allowbreak{}e\allowbreak{}-\allowbreak{}i\allowbreak{}r\allowbreak{}/\allowbreak{}1} live in Bijou, Geordi, Wesley, or a small shared package?
\item
  Should Bijou expose \texttt{b\allowbreak{}i\allowbreak{}j\allowbreak{}o\allowbreak{}u\allowbreak{}-\allowbreak{}l\allowbreak{}a\allowbreak{}y\allowbreak{}o\allowbreak{}u\allowbreak{}t\allowbreak{}/\allowbreak{}1} as a public artifact or keep it internal until the contract stabilizes?
\item
  Should Geordi consume \texttt{u\allowbreak{}i\allowbreak{}-\allowbreak{}s\allowbreak{}c\allowbreak{}e\allowbreak{}n\allowbreak{}e\allowbreak{}-\allowbreak{}i\allowbreak{}r\allowbreak{}/\allowbreak{}1} directly, or should Wesley emit \texttt{g\allowbreak{}e\allowbreak{}o\allowbreak{}r\allowbreak{}d\allowbreak{}i\allowbreak{}-\allowbreak{}i\allowbreak{}r\allowbreak{}/\allowbreak{}1} and \texttt{b\allowbreak{}i\allowbreak{}j\allowbreak{}o\allowbreak{}u\allowbreak{}-\allowbreak{}l\allowbreak{}a\allowbreak{}y\allowbreak{}o\allowbreak{}u\allowbreak{}t\allowbreak{}/\allowbreak{}1} as sibling artifacts?
\item
  How much action/focus metadata belongs in the shared IR versus the Bijou target profile?
\item
  Should debug JSONL files be produced by the renderer, by the lowering layer, or by both with different provenance phases?
\end{itemize}

\section{Recommended Next Step}\label{record-138-recommended-next-step}

Create a roadmap-backed issue for \texttt{u\allowbreak{}i\allowbreak{}-\allowbreak{}s\allowbreak{}c\allowbreak{}e\allowbreak{}n\allowbreak{}e\allowbreak{}-\allowbreak{}i\allowbreak{}r\allowbreak{}/\allowbreak{}1} and start with one tiny fixture scene. The fixture should prove that Bijou \texttt{S\allowbreak{}u\allowbreak{}r\allowbreak{}f\allowbreak{}a\allowbreak{}c\allowbreak{}e} is a render target for the shared IR, while Geordi remains the pixel-proof target.

